\section*{零基础读法与符号准备}
本册逐页保留原课件图像，再解释其数学含义。原图出现笔误时，在正文中明确纠正；截图与讲解的不同不代表遗漏。公式阅读先说中文：$\forall$ 是“对每一个”，$\exists$ 是“至少存在一个”，$\Rightarrow$ 是“能够推出”，$\Leftrightarrow$ 要证明两个方向。

\paragraph{最小工具箱}$x\in S$ 表示属于集合；$A\subseteq B$ 表示 $A$ 中每个元素都在 $B$ 中；$\bigcap_n A_n$ 要同时属于所有集合，$\bigcup_n A_n$ 只须属于至少一个。$\sup$ 是最小上界，$\inf$ 是最大下界，二者不保证属于原集合。绝对值 $|x-y|$ 表示一维距离。

\paragraph{怎样写证明}先抄清假设与目标，再把目标翻译成定义。证明“对所有”时先任取对象；证明“存在”时要构造对象并检查条件；反证法是暂时假定目标不成立并导出矛盾。否定“对每个 $\eps$ 存在 $K$”时，量词逐个反转，不能只在最后加一个“不”。

\paragraph{学习安排}每次阅读一个证明后，合上正文，回答：候选对象从哪里来？每个假设在哪里用到？删去一个假设有何反例？能回答这三问，再练习完整书写。原图中的英文术语可辅助对照，不要求先背英文定义。


\Slide{1}{实数论在课程中的地位}
本讲对应 Debreu《Theory of Value》第一章实数部分，并补上教材高度压缩的基础证明。实数完备性是后面极限、连续性、最优化和不动点存在性的基础。
标题页提供课程与教材信息，不包含需要证明的定理。学习本讲时要把“直觉知道”升级为“能从假设写出每一步”：例如知道 $1/n\to0$，还要能给定精度 $\eps$ 后选出一个有效的整数阈值。

\Slide{2}{章节路线图}
集合与映射提供语言，预序提供比较，实数提供大小和极限的基础。下一讲会从 $\R$ 推广到 $\R^N$，然后研究函数与对应的连续性、凸性和不动点。
本页是目录，因此不另造证明。建议先完成本讲确界、单调收敛、区间套、柯西、有限覆盖五组证明，再继续学紧集上的连续函数；没有完备性工具，后面的“存在最大值”很容易变成只会背结论。

\Slide{3}{实数的公理定义}
教材不从小数展开构造实数，而把实数规定为满足一组公理的对象。后面三页分别给出代数、顺序、完备性三部分。
“有限实数”意味着这里的 $\R$ 不含 $+\infty,-\infty$。写数列发散到 $+\infty$ 不等于它收敛到一个实数。遇到 $\sup S$ 时也要先问它是有限实数还是扩充实数约定；本讲确界定理要求相应方向有界，保证有限。

\Slide{4}{域公理：熟悉运算的严格起点}
加法、乘法都有交换律和结合律；加法单位元 $0$、乘法单位元 $1$ 不相等；每个数有负数，每个非零数有倒数；乘法对加法满足分配律。这些条件合起来说 $\R$ 是一个域。
\paragraph{基本推导}分配律给 $0x=(0+0)x=0x+0x$，在加法群中消去一项得 $0x=0$。若 $ab=0$ 且 $a\ne0$，乘以 $a^{-1}$ 得 $b=0$。这解释了“乘积为零则一个因子为零”为什么成立，也解释了不能除以零。
这些是证明时允许使用的基础。后面的代数术语会再次展开，不要求此时一次记住所有抽象名称。

\Slide{5}{全序与运算相容}
这里的 complete ordering 指全序：任意两个数都能比较，且满足自反、传递、反对称。它不是下一页的 Dedekind 完备性；“可比较”与“没有缺口”是不同概念。
\paragraph{相容性}不等式两边加同一个 $z$ 保持方向；非负数乘积非负。若 $a\le b,c\ge0$，因为 $b-a\ge0$，所以 $(b-a)c\ge0$，即 $ac\le bc$。若 $c<0$，先乘 $-c>0$ 再取相反数，得到 $ac\ge bc$。
\paragraph{考试提醒}乘不等式前检查乘数符号；未知符号的乘数不能直接保持方向。

\Slide{6}{完备性公理：有界集合存在边界}
非空上有界实数集合有最小上界，非空下有界实数集合有最大下界。上界是盖住全部元素的数，最小上界是所有这样数中最小的一个，未必属于原集合。
例如 $(0,1)$ 没有最大值，但有上确界 $1$。$\Q$ 也满足前面的域与序公理，却不能保证所有有界集合的边界仍在 $\Q$ 内：平方小于 $2$ 的正有理数集合的边界是 $\sqrt2$。
\paragraph{逻辑地位}这一页是公理，不从代数规则证明；后面会证明它与分割公理等价，并用它推出单调收敛、区间套等结论。

\Slide{7}{和、差、幂、倒数及唯一性}
$\sum_{n=1}^N x_n=x_1+\cdots+x_N$；$x^N$ 是 $N$ 个 $x$ 相乘，不是 $Nx$。减法定义为加负数，除法定义为乘倒数，且分母必须非零。
\paragraph{单位元唯一}若 $0'$ 也是加法单位元，则 $0=0+0'=0'$；乘法单位元同理。若 $u,v$ 都是 $x$ 的负数，则 $u=u+(x+v)=(u+x)+v=v$。若 $x\ne0$ 且 $u,v$ 都是其倒数，则 $u=u(xv)=(ux)v=v$。
\paragraph{证明 $(-1)x=-x$}由分配律，$x+(-1)x=(1-1)x=0x=0$，所以 $(-1)x$ 是 $x$ 的加法逆元；逆元唯一，因此等于 $-x$。课件省略的唯一性推理已经补全。

\Slide{8}{正、负、非负与非正}
$x\ge0$ 读作非负，允许为零；$x>0$ 才是严格正。非正与负数同理。$x\le y$ 等价于 $y-x\ge0$，这是运用不等式最方便的入口。
\paragraph{证明 $0<1$}任意非零数的平方严格为正：若 $a>0$，直接由序相容性；若 $a<0$，则 $-a>0$，且 $a^2=(-a)^2>0$。因为 $1\ne0$，有 $1=1^2>0$。
\paragraph{细节}非负乘积规则先给 $a^2\ge0$，而域中非零数的平方不为零，因此才能升级成严格正。写严格不等式时应说明为何等号不可能。

\Slide{9}{绝对值、符号和三角不等式}
最大元素 $m$ 必须满足 $m\in X$ 且每个 $x\in X$ 有 $x\le m$；若两个数都是最大元素，它们互相不大于对方，因此相等。最小元素同理。
$|x|=\max\{x,-x\}$，于是 $-|x|\le x\le|x|$。$x\ne0$ 时 $\operatorname{sign}x=x/|x|$，正数得 $1$，负数得 $-1$；本页没有给零定义该商。
\paragraph{三角不等式证明}将 $-|x|\le x\le|x|$ 与 $-|y|\le y\le|y|$ 相加，得到
\[-(|x|+|y|)\le x+y\le|x|+|y|.\]
一个数落在 $[-M,M]$ 中就等价于其绝对值不超过 $M$，因此 $|x+y|\le|x|+|y|$。后面所有误差估计都反复使用它。

\Slide{10}{确界与最值的区别}
上确界是最小上界；最大值是在集合内部达到的上确界。若 $a<b$，$[a,b)$ 有最小值 $a$、下确界 $a$、上确界 $b$，但无最大值，因为每个 $x<b$ 后面还有 $(x+b)/2$。
\paragraph{证明上确界是 $b$}首先所有 $x\in[a,b)$ 都满足 $x\le b$。若 $u<b$，当 $u<a$ 时 $a$ 已经超过它；当 $a\le u<b$ 时中点 $(u+b)/2$ 属于区间且超过 $u$。所以没有比 $b$ 更小的上界。
\paragraph{唯一性}若 $b,c$ 都是最小上界，则由彼此是上界得到 $b\le c$ 与 $c\le b$，因此 $b=c$。教材引用前一节的序理论，此处给出直接证明。

\Slide{11}{区间端点和长度}
有限端点为 $a\le b$ 的区间，其长度为 $b-a$；开闭端点不改变长度。$[a,b]$、$(a,b)$、$[a,b)$ 在 $a<b$ 时都有同样的上下确界。
\paragraph{解释}长度衡量跨度，端点是否被包含衡量边界归属，这两件事不同。区间套定理中“闭”控制极限点能否留在集合，“长度趋零”控制交点是否唯一。
\paragraph{边界情况}若 $a=b$，$[a,a]$ 是单点，长度零；$(a,a)$ 是空集，不能直接使用非空集合的确界定理。无界区间没有本讲意义下的有限长度。

\Slide{12}{整数嵌在实数中}
从 $0$ 反复加 $1$ 得 $1,2,\ldots$，反复减 $1$ 得 $-1,-2,\ldots$，合起来是 $\Z$。自然数在本课约定为 $\N=\{1,2,\ldots\}$，整数还包括 $0$ 与负整数。
\paragraph{可数性的具体排列}$0,1,-1,2,-2,3,-3,\ldots$ 把全部整数列成一个数列，因此 $\Z$ 可数。可数不表示有限，而表示能够按自然数编号。
\paragraph{常用整数性质}对任何实数 $t$ 存在整数 $m$ 使 $m-1\le t<m$。这由阿基米德性质与自然数的最小元性质得到，在证明两个实数之间存在有理数时会用到。

\Slide{13}{有理数可数，但不完备}
每个有理数都可写成 $p/q$，其中 $p\in\Z,q\in\N$。按 $|p|+q=1,2,3,\ldots$ 分层列出所有这样的整数对；每层有限，所有层可以顺序列举。把各对映成 $p/q$ 并删除重复项，便枚举了 $\Q$，证明其可数。
\paragraph{为什么没有有理平方根 $\sqrt2$}若 $(p/q)^2=2$ 且分数已约成互素，得 $p^2=2q^2$，所以 $p$ 偶，写 $p=2k$；代回得 $q^2=2k^2$，所以 $q$ 也偶，与互素矛盾。
\paragraph{不完备要再补一步}令 $S=\{r\in\Q:r^2<2\}$。若其有理上确界为 $c$，则 $c\ge1$。若 $c^2<2$，取充分小的正有理数 $h$，使 $h<1$ 且 $(2c+1)h<2-c^2$，则 $(c+h)^2<2$，与上界矛盾。若 $c^2>2$，取 $0<h<c$ 且 $2ch<c^2-2$，则 $(c-h)^2>2$；所有 $r\in S$ 都小于 $c-h$，产生更小上界，矛盾。故必须 $c^2=2$，而刚证无有理解。因此 $S$ 在 $\Q$ 中没有上确界。

\Slide{14}{实数不可数与有理数稠密}
\paragraph{先证阿基米德性质}$\N$ 不可能上有界：若 $s=\sup\N$，存在 $n>s-1$，则 $n+1>s$，矛盾。因此对任意实数 $t$ 都有整数 $n>t$。
\paragraph{有理数稠密的证明}给 $x<y$，取正整数 $n$ 使 $n(y-x)>1$。取整数 $m$ 使 $m-1\le nx<m$，则 $nx<m\le nx+1<ny$。除以正数 $n$ 得 $x<m/n<y$。任意短的非空实数区间都包含有理数，这称为稠密。
\paragraph{实数不可数的证明}反设 $[0,1]$ 可列为 $r_1,r_2,\ldots$。从 $I_0=[0,1]$ 开始，每次把当前闭区间的左三分之一与右三分之一作为两个不相交候选区间，至少一个不含 $r_n$，选它作 $I_n$。则 $I_n$ 非空、闭、嵌套，长度 $3^{-n}$。后文证明的区间套定理给交点 $a$，而 $a\not=r_n$ 对每个 $n$ 成立，矛盾。该证明依赖区间套定理，可在读完后面证明后回来看；它没有依赖不可数性来证明区间套，故不循环。
\paragraph{区分}稠密描述“无处不在”，可数描述“能否编号”；有理数同时稠密且可数，二者不矛盾。

\Slide{15}{非负数的整数次根：存在与唯一}
给 $x\ge0$、整数 $N\ge1$。$x=0$ 时根为 $0$，$N=1$ 时根为 $x$。以下设 $x>0$。令 $S=\{t\ge0:t^N\le x\}$。它非空（含 $0$），且有上界，例如 $\max(1,x)+1$；取 $y=\sup S$。
\paragraph{关键估计}对 $a,b\in[0,M]$，利用因式分解得到
\[|a^N-b^N|=|a-b|\left|\sum_{j=0}^{N-1}a^{N-1-j}b^j\right|\le NM^{N-1}|a-b|.\]
\paragraph{排除 $y^N<x$}取足够小 $h>0$，使 $h<1$ 且 $N(y+1)^{N-1}h<x-y^N$，上述估计给 $(y+h)^N<x$，违背 $y$ 为上界。
\paragraph{排除 $y^N>x$}先有 $y>0$。取足够小 $0<h<y$，使 $Ny^{N-1}h<y^N-x$，则 $(y-h)^N>x$。由于非负数上 $t\mapsto t^N$ 严格递增，所有 $t\in S$ 均小于 $y-h$，这产生更小上界，矛盾。因此 $y^N=x$。
\paragraph{唯一性}若 $0\le a<b$，则 $b^N-a^N=(b-a)\sum_{j=0}^{N-1}b^{N-1-j}a^j>0$。所以两个不同非负数不可能有相同的 $N$ 次幂。这个证明直接用确界与代数，不预先借用尚未证明的连续性定理。

\Slide{16}{数列收敛定义的量词顺序}
$x_q\to x_0$ 的精确定义是：$\forall\eps>0\ \exists K\in\N\ \forall q>K$，有 $|x_q-x_0|<\eps$。先由对方提出精度，再由你选阈值；阈值以后所有项都要过关，不能只找某一项。
\paragraph{示范}证明 $1/q\to0$。给 $\eps>0$，由阿基米德性质取整数 $K>1/\eps$。当 $q>K$，$|1/q-0|=1/q<\eps$。
\paragraph{如何否定}不收敛到 $x_0$ 的意思是存在一个固定 $\eps_0>0$，无论给多大的 $K$，仍有某个 $q>K$ 满足 $|x_q-x_0|\ge\eps_0$。注意固定精度在最外层；它不是随着 $K$ 缩小的误差。

\Slide{17}{极限唯一性与三个例子}
\paragraph{极限唯一性证明}若 $x_q\to a$ 又 $x_q\to b$ 且 $a\ne b$，取 $\eps=|a-b|/3$。充分大的同一个 $q$ 同时满足 $|x_q-a|<\eps$、$|x_q-b|<\eps$。三角不等式给 $|a-b|<2|a-b|/3$，矛盾。
\paragraph{交替但收敛}$|(-1)^q/q|=1/q$，给定精度后取 $q>1/\eps$，故极限为零。上下摆动不妨碍收敛，只要幅度趋零。
\paragraph{下确界不等于最小值}$X=\{1/q:q\in\N\}$ 的每个元素为正，所以 $0$ 是下界；对任意 $\eps>0$ 有 $1/q<\eps$，故 $\inf X=0$。但 $0\notin X$，且每个 $1/q$ 后都有更小的 $1/(q+1)$，所以没有最小值。
\paragraph{不收敛}$(-1)^q$ 的偶数子列恒为 $1$、奇数子列恒为 $-1$，不可能有共同极限；$q$ 无界，而收敛数列必有界（尾部距极限小于 $1$，前面有限项再取最大值），所以也不收敛。

\Slide{18}{数轴把绝对值解释为距离}
把 $0$ 与 $1$ 固定后，正数在右、负数在左，$|x|$ 是离原点的距离，$|x-y|$ 是两点间距离。条件 $|x-a|<\eps$ 等价于 $a-\eps<x<a+\eps$，也就是进入以 $a$ 为中心的一个小邻域。
这幅几何图景解释极限定义：不论给目标点画多窄的邻域，数列后段都会全部留在里面。数学证明仍需写出阈值或估计；一张图可以帮助猜测，不能替代对所有精度的论证。

\Slide{19}{四个数系与符号约定}
本书固定使用 $\N$ 表示正整数，$\Z$ 表示整数，$\Q$ 表示有理数，$\R$ 表示实数，且 $\N\subset\Z\subset\Q\subset\R$。同一个符号在其他书中可能有不同约定，例如有些书把 $0$ 算入自然数。
\paragraph{本段小结}你应已经能证明：负数与倒数唯一、三角不等式、确界唯一、有理数可数与稠密、极限唯一。后面的课程补充会系统建立完备性的诸多等价表现，并把刚才用过的区间套证明补齐。

\Slide{20}{从自然数逐步构造实数}
$\N=\{1,2,\ldots\}$ 用来计数；加入 $0$ 与负数得到 $\Z$；允许非零整数作分母得到 $\Q$。但 $\Q$ 中没有平方等于 $2$ 的数，说明只会四则运算还不足以容纳所有极限。
\paragraph{构造的直觉}用 $1,1.4,1.41,1.414,\ldots$ 逐渐逼近 $\sqrt2$。每一项都是有理数，目标却不在 $\Q$ 中。柯西序列的方法把这样的“越来越稳定的逼近”作为新数。
\paragraph{准确补充}严格地说，实数不是所有有理柯西序列本身，而是这些序列的\emph{等价类}：$(a_n)\sim(b_n)$ 当 $a_n-b_n\to0$。不同逼近过程可能代表同一个数。Dedekind 分割则用目标左侧和右侧的有理数来定位目标；本课使用性质而不要求从头建设整个数系。

\Slide{21}{实数是完备有序域}
“域”保证加减乘除（除数非零）合法；“有序”保证能比较大小且运算保持相容；“完备”保证有界的逼近不会落进一个缺口。三个层次缺一不可。
\paragraph{为什么必须强调完备}有理数也是有序域，却不是完备有序域。因而证明中一旦说“取 $\sup S$”或“这个柯西序列有极限”，实际上就用了比四则运算更强的性质。
\paragraph{证明意识}公理是起点，不要求在同一体系内再证明。选择确界公理后，单调收敛、闭区间套等才是要从它推出的定理；不能互相引用绕成一圈。

\Slide{22}{加法群：能加、能撤销}
封闭性要求 $a+b$ 仍在集合里；结合律只改变括号；单位元 $0$ 加上去不改变数；逆元 $-a$ 撤销加上 $a$ 的效果。交换律 $a+b=b+a$ 是另外一个条件，满足它才叫交换群。
\paragraph{两个小证明}若 $e,e'$ 都是单位元，则 $e=e+e'=e'$，故单位元唯一。若 $b,c$ 都是 $a$ 的加法逆元，则
\[b=b+0=b+(a+c)=(b+a)+c=0+c=c.\]
每个等号依次使用单位元、逆元、结合律、逆元、单位元。这就是“由公理证明”的写法。
\paragraph{例子与反例}$\Z$ 对加法是群；$\N$ 对加法不是群，因为缺 $0$ 和负数。判断群时必须同时说明集合与运算。

\Slide{23}{乘法群：为什么排除零}
乘法单位元是 $1$，乘法逆元是 $a^{-1}=1/a$。$\R\setminus\{0\}$ 对普通乘法是交换群，$\R$ 对普通乘法不是群，因为 $0b=0$ 不可能等于 $1$。
\paragraph{逆元唯一}若 $ab=ba=1$ 且 $ac=ca=1$，则 $b=b(ac)=(ba)c=c$。这里不能把“逆元”与“负数”混淆：$2$ 的加法逆元是 $-2$，乘法逆元是 $1/2$。
\paragraph{自检}$\Z\setminus\{0\}$ 对乘法不是群，因为 $2$ 没有整数逆元；正实数对乘法是群，因为正数的乘积及倒数仍为正数。

\Slide{24}{一般群与消去律}
本页把加法、乘法统一写成运算 $\cdot$。符号不意味着一定是普通乘法；例如置换的复合也是群运算。$e$ 是单位元，$a^{-1}$ 是相对于该运算的逆元。
\paragraph{左消去律的证明}若 $a\cdot b=a\cdot c$，两边左乘 $a^{-1}$，由结合律得
\[(a^{-1}\cdot a)\cdot b=(a^{-1}\cdot a)\cdot c,\qquad e\cdot b=e\cdot c,\]
因此 $b=c$。右消去同理，但应从右侧乘逆元。在非交换群中不能任意调换因子。
\paragraph{层次}封闭性、结合律、单位元、逆元是定义；唯一性和消去律是由定义推出的性质。考试写证明时要分清两者。

\Slide{25}{环：把加法和乘法放在一起}
本课采用有乘法单位元的环定义：加法构成交换群；乘法封闭、有结合律和单位元；两种运算满足左右分配律。一般环的乘法不必交换，也不要求非零元素都有倒数。
\paragraph{基本推论}由 $a(0+0)=a0+a0$ 及 $0+0=0$，得 $a0=a0+a0$。两边加上 $-(a0)$ 得 $a0=0$。同理 $0a=0$。
\paragraph{例子}$\Z$ 是环，但 $2$ 在 $\Z$ 中不可逆；实方阵集合也是环，矩阵乘法一般不交换。不要把这里表示任意环的字母 $R$ 自动读成实数集 $\R$。

\Slide{26}{交换环比一般环多了什么}
本页比上一页增加乘法交换律 $ab=ba$，而加法交换律在环定义中已经有了。整数、有理数、实数、多项式都提供熟悉的例子。
\paragraph{不要误用约分}在一般交换环里，$ab=ac$ 且 $a\ne0$ 未必推出 $b=c$。模 $6$ 的余数环中 $2\cdot1=2\cdot4$，但 $1\ne4$。失败来自 $2\cdot3=0$ 这样的零因子。
\paragraph{考试检查}写“除以 $a$”前，应确认 $a$ 有乘法逆元；只有非零还不一定足够。下一页的域正是让每个非零元素都可逆的结构。

\Slide{27}{域：合法四则运算的基础}
域是满足 $0\ne1$ 的交换环，并且每个非零元素有乘法逆元。因此 $a/b$ 的含义是 $a\cdot b^{-1}$，条件 $b\ne0$ 不能省略。
\paragraph{无零因子的证明}设 $ab=0$。若 $a=0$，结论成立；若 $a\ne0$，左乘 $a^{-1}$ 得 $b=a^{-1}0=0$。因此域中 $ab=0$ 必有一个因子为零。这一性质使许多代数约分合理。
\paragraph{辨认}$\Q,\R,\mathbb C$ 都是域，$\Z$ 不是域。本页把三组条件称作 ring axioms 属于不够精确的表述；含逆元条件后已是域公理。末行变量所属集合应为 $F$。

\Slide{28}{整环、除环和域}
整环是无零因子的含幺交换环；除环要求每个非零元素可逆，但乘法不必交换；交换的除环就是域。因而“整环”不等于“域”，$\Z$ 是直接反例。
\paragraph{关系要这样读}$\Z\subseteq\Q\subseteq\R\subseteq\mathbb C$ 是集合包含，并且普通四则运算在这些集合间相容。所有列出的数系都是环，只有后面三个是域。
\paragraph{小练习及答案}证明整环中非零元素可消去。若 $ab=ac$，则 $a(b-c)=0$；由 $a\ne0$ 和无零因子得 $b-c=0$。注意这里证明了消去律，却没有创造出一个不存在的整数 $1/2$。

\Slide{29}{向量空间中的两类对象}
向量空间包含向量 $v\in V$，其标量来自域 $F$。向量相加是 $V\times V\to V$；数乘是 $F\times V\to V$。例如 $V=\R^2,F=\R$，$2(1,3)=(2,6)$。
\paragraph{最易混淆处}数乘的第一个输入是数，第二个是向量；通常没有定义两个向量相乘仍为向量。经济学中的商品束 $(x_1,\ldots,x_N)$ 就可放在 $\R^N$ 中。
\paragraph{例子}所有连续实函数也构成向量空间，运算逐点定义：$(f+g)(t)=f(t)+g(t)$。因此“向量”不一定是画在平面上的箭头，关键是满足下页的运算法则。

\Slide{30}{向量空间八条公理如何使用}
前四条说明向量加法是交换群；后四条说明数乘与域运算配合。所有公理都要注意变量类型：$v,v_1,v_2\in V$，$a,a_1,a_2\in F$。
\paragraph{推导 $0v=0_V$}利用分配律，$0v=(0+0)v=0v+0v$，在向量加法群中消去一个 $0v$ 即得。类似地，$a0_V=a(0_V+0_V)=a0_V+a0_V$，所以 $a0_V=0_V$。又因为 $v+(-1)v=(1-1)v=0_V$，故 $(-1)v=-v$。
\paragraph{原页勘误}第(1)条把向量写成属于 $F$，应为属于 $V$。两个“零”分别属于域和向量空间，证明时写 $0_V$ 可避免混淆。

\Slide{31}{用正数集合定义大小}
$P$ 是严格正数集。三歧性说一个数恰好为零、为正或其相反数为正；正数之和、积仍为正。定义 $a<b$ 等价于 $b-a\in P$，把大小关系转成正数运算。
\paragraph{传递性的证明}若 $a<b$ 且 $b<c$，则 $b-a,c-b\in P$。正数相加仍正，故 $(b-a)+(c-b)=c-a\in P$，即 $a<c$。
\paragraph{不等式法则}由 $(b+z)-(a+z)=b-a$，同时加 $z$ 不改变大小。若 $c>0$，则 $bc-ac=c(b-a)>0$；若 $c<0$，乘上 $c$ 会反向。这里 $\R_+=[0,\infty)$，$\R_{++}=(0,\infty)$，须服从老师约定。

\Slide{32}{有序域与复数的区别}
有序域是域结构和满足相容条件的全序合在一起。$\Q$ 与 $\R$ 都是有序域，但只有 $\R$ 完备。
\paragraph{为什么复数不能这样排序}在有序域中，任意非零 $a$ 的平方为正：$a>0$ 时直接用正数乘法；$a<0$ 时 $a^2=(-a)^2>0$。因此 $1=1^2>0$，而 $i^2=-1$ 不可能也为正。故不能给 $\mathbb C$ 配上与域运算相容的实数式大小关系。
\paragraph{符号说明}课件的 $0_\oplus$ 是加法单位元，$0_\otimes$ 表示乘法单位元，后者通常写作 $1$；这不表示实数的 $0$ 与 $1$ 相等。

\Slide{33}{Dedekind 分割：数轴没有断口}
把全部实数分为非空的左部 $A$ 和右部 $B$，每个左部数都小于每个右部数。完备性说存在唯一分界点 $c$：小于它的都在 $A$，大于它的都在 $B$，而 $c$ 自己必须分到一边。
\paragraph{唯一性证明}若 $c<d$ 都能作分界点，取中点 $m=(c+d)/2$。由 $m>c$ 得 $m\in B$，由 $m<d$ 得 $m\in A$，矛盾。因为原条件使 $A\cap B=\varnothing$。
\paragraph{为什么两部分非空}若 $A=\R,B=\varnothing$，就没有有限分界点。若在 $\Q$ 中用平方小于或大于 $2$ 的条件切割，缺失的 $\sqrt2$ 正说明相同结论不适用于 $\Q$。

\Slide{34}{分界点究竟属于哪边}
本页把分割公理写成“左边有最大值，或右边有最小值”。若 $c\in A$，则 $c=\max A$、$A=(-\infty,c]$；若 $c\in B$，则 $c=\min B$、$B=[c,\infty)$。
\paragraph{为什么不能两边同时有端点}若 $a=\max A$ 且 $b=\min B$，分割条件给 $a<b$，但中点 $(a+b)/2$ 既不能在 $A$ 也不能在 $B$，与并集是整个实数轴矛盾。
$A^c$ 表示相对于 $\R$ 的补集，即 $\R\setminus A$。这页没有引入另一种完备性，只是上一页的等价说法。

\Slide{35}{上界、下界、上确界、下确界}
对集合 $S$，上界 $u$ 要满足“每一个 $x\in S$ 都有 $x\le u$”；上确界 $b$ 除了是上界，还必须小于等于任何其他上界。下界与下确界把不等号反过来。
\paragraph{例子}$S=(0,1)$ 的上界是 $[1,\infty)$，最小的上界是 $1$，故 $\sup S=1$；其下界是 $(-\infty,0]$，故 $\inf S=0$。但 $S$ 没有最大、最小元素，因为 $0,1$ 不在其中。
\paragraph{唯一性}若 $b,c$ 都是上确界，由 $c$ 是上界得 $b\le c$，交换角色得 $c\le b$，所以 $b=c$。证明一个候选数是确界时，必须检查“界”与“最优的界”两部分。

\Slide{36}{确界存在定理的假设}
非空且有上界的实数集合有有限实数上确界；非空且有下界的实数集合有有限实数下确界。“非空”“有界的方向”“实数”都是实质条件。
\paragraph{由上确界推出下确界}令 $-S=\{-x:x\in S\}$。若 $S$ 非空且有下界，则 $-S$ 非空且有上界，设 $b=\sup(-S)$。对每个 $x\in S$，$-x\le b$，故 $x\ge-b$；若 $d$ 是 $S$ 的下界，则 $-d$ 是 $-S$ 的上界，故 $b\le-d$，即 $d\le-b$。所以 $\inf S=-\sup(-S)$。
\paragraph{考试提醒}确界不一定在集合中。“存在确界”不能改写成“存在最大值”。

\Slide{37}{由分割公理证明上确界存在}
\paragraph{构造}令 $B$ 为 $S$ 的全部上界，$A=\R\setminus B$。有上界使 $B\ne\varnothing$；选 $s\in S$，$s-1$ 不是上界，故 $A\ne\varnothing$。
\paragraph{验证分割}若 $a\in A,b\in B$ 而 $a\ge b$，则所有 $s\in S$ 满足 $s\le b\le a$，使 $a$ 也成上界，矛盾。因此 $a<b$，分割公理给分界点 $c$。
\paragraph{证明端点属于上界侧}若 $c\notin B$，有 $s\in S$ 满足 $s>c$。中点 $m=(c+s)/2$ 大于 $c$，应属于 $B$；但 $m<s$，又不是上界，矛盾。故 $c\in B$。所有 $b\in B$ 都有 $c\le b$，所以 $c=\min B=\sup S$。
关键技巧：用“全部上界”形成新集合，不是直接把原集合随意切开。

\Slide{38}{由上确界存在证明分割公理}
设 $A,B$ 是一个分割。$A$ 非空且任取 $b\in B$ 都是其上界，故可取 $c=\sup A$。
\paragraph{验证分界性质}若 $x<c$ 而 $x\in B$，则 $x$ 是 $A$ 的上界，违背 $c$ 是最小上界，所以 $x\in A$。若 $x>c$，由于 $c$ 是 $A$ 的上界，$x$ 不可能在 $A$，只能在 $B$。
\paragraph{唯一性补全}两个不同分界点之间的中点会被同时分入 $A$ 和 $B$，矛盾。故确界公理能推出完整的分割公理。考试中“等价”必须写出两个方向，本页与前页缺一不可。

\Slide{39}{由分割公理证明下确界存在}
令 $A$ 为 $S$ 的全部下界，$B=\R\setminus A$。由下有界，$A\ne\varnothing$；取 $s\in S$，$s+1$ 不是下界，故 $B\ne\varnothing$。若 $a\in A,b\in B$ 且 $a\ge b$，则 $b\le a\le s$ 对所有 $s\in S$ 成立，矛盾，故 $a<b$。
设分界点为 $c$。如果 $c\in B$，则有 $s\in S$ 满足 $s<c$。中点 $m=(s+c)/2$ 小于 $c$，应在 $A$，但 $m>s$ 又不可能是下界。因此 $c\in A$，且所有下界 $a\le c$，故 $c=\inf S$。
\paragraph{记忆方法}上确界证明中选择“全部上界”的最小者；这里选择“全部下界”的最大者。中点总是用来挤出与分界方向冲突的数。

\Slide{40}{由下确界存在恢复分割}
$B$ 非空且任一 $a\in A$ 是其下界，取 $c=\inf B$。若 $x<c$，$x$ 不可能属于 $B$，否则违反 $c$ 是下界，所以 $x\in A$。若 $x>c$ 而 $x\in A$，则 $x$ 是 $B$ 的一个\emph{大于} $c$ 的下界，违背 $c$ 为最大下界，因此 $x\in B$。
\paragraph{原页勘误}末段英文写成 ``less than the greatest lower bound''，应为 ``greater than''。证明中的矛盾恰好来自找到一个更大的下界。
\paragraph{练习}把每个“不能”改写成反证的完整一句。这样能检查究竟使用了“$c$ 是下界”还是“$c$ 是最大下界”。

\Slide{41}{确界的逼近刻画}
$b=\sup S$ 当且仅当：所有 $x\in S$ 满足 $x\le b$，而且每给 $\eps>0$，都能找到 $x\in S$ 使 $b-\eps<x\le b$。下确界类似：$\ell\le x<\ell+\eps$。
\paragraph{直觉}集合可能碰不到边界，但一定能从内部任意接近边界。$\eps$ 是任意精度要求；找到的 $x$ 可以随 $\eps$ 改变，不要求同一个 $x$ 满足所有精度。
\paragraph{量词陷阱}“每个 $\eps$ 都有一个 $x$”写作 $\forall\eps>0\ \exists x\in S$；若颠倒成 $\exists x\ \forall\eps>0$，便强迫 $x=b$，这就成了最大值，强得多。后两页分别证明这两个等价刻画。

\Slide{42}{上确界逼近刻画的双向证明}
\paragraph{必要性}设 $b=\sup S$。它当然是上界。如果某个 $\eps>0$ 下找不到 $x>b-\eps$，那么所有 $x\in S$ 都满足 $x\le b-\eps$，说明 $b-\eps$ 是更小上界，矛盾。
\paragraph{充分性}设 $b$ 是上界且有逼近性质。任取上界 $u$，要证 $b\le u$。若 $u<b$，令 $\eps=b-u>0$。逼近性质给 $x\in S$ 满足 $x>b-\eps=u$，违反 $u$ 是上界。因此 $b\le u$。
\paragraph{如何想到选精度}想让 $b-\eps$ 恰好落在反证假设中的 $u$，所以解方程得到 $\eps=b-u$。证明中不是“随便取一个小数”，而是围绕目标构造。

\Slide{43}{下确界逼近刻画的双向证明}
\paragraph{必要性}若 $\ell=\inf S$，它是下界。若某个 $\eps>0$ 使所有 $x\in S$ 都满足 $x\ge\ell+\eps$，那么 $\ell+\eps$ 是更大的下界，矛盾。
\paragraph{充分性}若 $\ell$ 是下界并能从上方逼近，任取下界 $d$。假设 $d>\ell$，令 $\eps=d-\ell$，就有 $x\in S$ 满足 $x<\ell+\eps=d$，违背下界性质。故每个下界 $d\le\ell$。
\paragraph{一眼识别使用时机}需要构造 $x_n\in S$ 逼近下确界时，逐次取 $\eps=1/n$，得到 $\ell\le x_n<\ell+1/n$，于是 $x_n\to\ell$。存在逼近数列不意味着下确界属于 $S$。

\Slide{44}{单调与有界是两种不同限制}
数列是按次序排列的无限列表 $(x_1,x_2,\ldots)$。递增在本课指 $x_n\le x_{n+1}$，允许相等；严格递增才要求 $<$。有界指存在一个固定 $M$，使每个 $n$ 都有 $|x_n|\le M$。
\paragraph{例子}$1-1/n$ 递增且上有界；$1/n$ 递减且下有界；$n$ 递增但上无界；$(-1)^n$ 有界但不单调。后两例说明只满足一半条件不能保证有限极限。
\paragraph{预备定义}写 $x_n\to a$ 的完整意思是：对每个 $\eps>0$，存在整数 $K$，使所有 $n\ge K$ 满足 $|x_n-a|<\eps$。$K$ 可以依赖 $\eps$，但不能依赖随后任取的 $n$。

\Slide{45}{单调有界收敛定理的准确版本}
递增且有上界的数列收敛到全体项的上确界；递减且有下界的数列收敛到全体项的下确界。无需再额外假设递增数列下有界，因为第一项就是下界；递减情形同理。
\paragraph{证明策略}先用确界公理产生候选极限，再用确界逼近性质找到“至少一个”接近候选值的项，最后用单调性把“至少一个”升级成“此后全部”。这三个步骤是本定理的核心。
\paragraph{避免倒因为果}不能先假定极限存在再算极限。本定理正是用来证明存在的；例如递推数列先证明单调有界后，才可在递推式中取极限。

\Slide{46}{递增情形：逐句拆解证明}
设 $a=\sup\{x_n:n\in\N\}$，存在性来自非空和上有界。由上界定义，对所有 $n$ 有 $x_n\le a$。
任给 $\eps>0$，由于 $a-\eps<a$，它不可能仍是上界。因此存在某个 $K$，使 $x_K>a-\eps$。单调性告诉我们：只要 $n\ge K$，就有 $x_n\ge x_K>a-\eps$。合起来：
\[a-\eps<x_K\le x_n\le a,\qquad |x_n-a|=a-x_n<\eps.\]
这正是极限定义，故 $x_n\to a$。
\paragraph{考试易漏}确界性质只给一个接近上界的项；如果没有单调性，后面的项可能又跑远。不能把这一步省成“由确界定义显然收敛”。

\Slide{47}{递减情形与课件笔误}
\paragraph{原页勘误}证明首段的 ``supremum'' 应改为 ``infimum''；后面的 ``increasing'' 应改为 ``decreasing''。
\paragraph{完整证明}令 $a=\inf\{x_n:n\in\N\}$。非空且下有界保证存在，并且 $a\le x_n$。任给 $\eps>0$，由下确界逼近性质，存在 $K$ 使 $x_K<a+\eps$。当 $n\ge K$，递减性给 $x_n\le x_K$，因此
\[a\le x_n\le x_K<a+\eps,\qquad |x_n-a|=x_n-a<\eps.\]
故 $x_n\to a$。
\paragraph{自检}以 $x_n=1/n$ 检查：其极限是下确界 $0$，绝不是上确界 $1$。代入简单例子常能发现符号错误。

\Slide{48}{闭区间套的四种表述}
闭区间套是非空闭区间 $I_n=[d_n,u_n]$ 满足 $I_{n+1}\subseteq I_n$。于是左端点不减、右端点不增：$d_n\le d_{n+1}\le u_{n+1}\le u_n$。
结论先是交集非空，进一步精确为
\[\bigcap_{n=1}^{\infty}I_n=[D,U],\qquad D=\sup_n d_n,\quad U=\inf_n u_n.\]
若区间长度 $u_n-d_n\to0$，则 $D=U$，交集恰有一个点。
\paragraph{假设不能删除}开区间 $(0,1/n)$ 嵌套但交集为空；闭区间 $[n,\infty)$ 嵌套但不有界，交集为空；恒定区间 $[0,1]$ 说明不缩短到零就不保证唯一。

\Slide{49}{先把左端点的上确界放入每个区间}
所有 $d_n\le u_1$，因此 $D=\sup_n d_n$ 存在，并且 $d_m\le D$。还要证明 $D\le u_m$，才可说 $D\in I_m$。
\paragraph{关键交叉比较}固定任意 $m$，对任意 $n$：若 $n\le m$，则 $d_n\le d_m\le u_m$；若 $n>m$，则 $d_n\le u_n\le u_m$。所以 $u_m$ 是\emph{全部}左端点的上界，故 $D\le u_m$。
于是每个 $m$ 都有 $d_m\le D\le u_m$，即 $D\in\bigcap_m I_m$。
\paragraph{原页勘误}末尾英文 union 应为 intersection（交集）。只比较同一指标 $d_n\le u_n$ 不够，必须证明任意左右端点的交叉不等式。

\Slide{50}{再把右端点的下确界放入每个区间}
\paragraph{纠正后证明}所有 $u_n\ge d_1$，所以右端点集合\emph{下有界}，$U=\inf_n u_n$ 存在，且 $U\le u_m$。上一页的交叉比较给固定 $m$ 下，所有 $n$ 满足 $d_m\le u_n$；因此 $d_m$ 是右端点集合的下界，故 $d_m\le U$。所以 $d_m\le U\le u_m$，$U$ 属于每个区间。
\paragraph{原页勘误}bounded above 应改为 bounded below；后面的 $\sup B$ 应为 $\inf B$；$l_n$、$l_{n'}$ 应统一为左端点 $d_n,d_{n'}$；关键结论是 $d_m\le\inf B=U$，不是 $u_m\le U$。
本证明只用下确界定义，没有偷偷假设交集非空。

\Slide{51}{为什么交集恰好是一个闭区间}
要证两个集合相等，标准做法是证明两个包含关系。
\paragraph{第一方向}若 $x\in[D,U]$，每个 $n$ 都有 $d_n\le D\le x\le U\le u_n$，所以 $x\in I_n$，因此 $x$ 属于交集。
\paragraph{第二方向}若 $x\in\bigcap_n I_n$，则每个 $n$ 都有 $d_n\le x$，所以 $x$ 是全体左端点的上界，从而 $D\le x$。同理 $x\le u_n$ 对所有 $n$ 成立，使 $x$ 成为右端点集合的下界，所以 $x\le U$。故 $x\in[D,U]$。
这里的第二方向用确界定义直接完成，比课件的两次反证更短，但逻辑相同。$D\le U$ 已由左端点的上确界属于所有区间得到。

\Slide{52}{长度趋零如何产生唯一交点}
已有 $d_n\le D\le U\le u_n$，两端相减可得
\[0\le U-D\le u_n-d_n.\]
右边趋于 $0$，所以固定非负数 $U-D$ 只能为 $0$，即 $D=U=:a$。交集 $[D,U]=\{a\}$，且单调收敛定理给 $d_n\to a,u_n\to a$。
\paragraph{不引用夹逼定理也能证明}若 $U-D=c>0$，由长度趋零，某个 $n$ 有 $u_n-d_n<c$，与 $c\le u_n-d_n$ 矛盾。
\paragraph{“唯一”的证明}若 $a,b$ 都在每个区间，则 $|a-b|\le u_n-d_n$ 对所有 $n$ 成立，同样推出 $a=b$。长度条件在这一步才真正用到。

\Slide{53}{Bolzano 定理：二分法抽子列}
把全部项放进 $[D,U]$。若 $D=U$，数列常值，直接收敛。否则把区间二等分，至少一半含无穷多个\emph{项的指标}；否则两半加起来仅含有限项。保留这一半，重复二分，得到长度 $(U-D)/2^k$ 的闭区间套 $I_k$。
每个 $I_k$ 含无穷多个指标，所以可以逐次选 $n_k>n_{k-1}$ 且 $x_{n_k}\in I_k$。区间套定理给唯一交点 $a$。因为 $a,x_{n_k}$ 都在 $I_k$，
\[|x_{n_k}-a|\le |I_k|=(U-D)2^{-k}\longrightarrow0.\]
故选出的子列收敛。
\paragraph{细节}原页的严格 $<$ 一般应改为 $\le$，不影响极限。常值数列只有一个不同取值，却有无穷多个项；所以这里不能误说“无穷多个不同点”。

\Slide{54}{有界无限集合必有聚点：构造阶段}
聚点 $a$ 的定义是：每个 $\delta>0$ 都能找到 $x\in E$ 满足 $0<|x-a|<\delta$。允许 $a$ 不属于 $E$，但找出的点必须不同于 $a$。
把有界无限集合 $E$ 放入 $[-M,M]$，取 $M>0$。二分后至少一个闭半区间仍含 $E$ 的无穷多个\emph{不同点}，否则两边有限，其并也有限。反复保留无限的一半，得到 $I_n$ 嵌套，长度 $2M/2^n$，且 $I_n\cap E$ 无限。
\paragraph{与上一页的区别}数列版本保留无限指标；集合版本保留无限不同点。这一区别保证下一页能排除交点本身，真正证明“聚点”。原页 $R_0$ 应读作 $I_0$。

\Slide{55}{有界无限集合必有聚点：验证阶段}
闭区间套给唯一交点 $a$。任给 $\delta>0$，选择 $n$ 使 $|I_n|<\delta$。对每个 $x\in I_n$，因 $a\in I_n$，有 $|x-a|\le|I_n|<\delta$，因此 $I_n\subset(a-\delta,a+\delta)$。
又因为 $E\cap I_n$ 有无穷多个不同点，删去可能存在的点 $a$ 后仍至少有一个 $x$。这个 $x$ 满足 $x\in E$ 且 $0<|x-a|<\delta$。由于 $\delta$ 任意，$a$ 是聚点。
\paragraph{课件省略的关键句}只说邻域与集合相交尚不足以证明聚点，孤立点的邻域也会与集合相交。必须说“有不同于 $a$ 的点”，无限性正是用来保证这一点。

\Slide{56}{柯西列：不用预先知道极限的收敛标准}
柯西条件是
\[\forall\eps>0\ \exists K\ \forall m,n\ge K:\quad |x_n-x_m|<\eps.\]
收敛比较“项与目标”；柯西比较“后段任意两项”。因而当我们还不知道目标在哪里时，柯西条件特别有用。
\paragraph{不要只比较相邻项}$|x_{n+1}-x_n|\to0$ 比柯西弱得多。例如 $x_n=\sqrt n$ 的相邻差为 $1/(\sqrt{n+1}+\sqrt n)\to0$，但数列无界，不收敛。柯西要求任意远的两个后段项也很接近。
\paragraph{完备性的位置}收敛必为柯西在任何度量空间成立；柯西必收敛到空间内某点，是 $\R$ 的完备性，不适用于 $\Q$。

\Slide{57}{收敛推出柯西：把误差分成两半}
设 $x_n\to a$。给定最终希望得到的 $\eps>0$，按收敛定义，对精度 $\eps/2$ 选 $K$，使 $n\ge K$ 时 $|x_n-a|<\eps/2$。
若 $m,n\ge K$，利用三角不等式：
\[|x_n-x_m|=|(x_n-a)+(a-x_m)|\le|x_n-a|+|x_m-a|<\eps.\]
于是满足柯西定义。
\paragraph{为什么课件写 $2\eps$ 也能成立}因为 $\eps$ 是任意正数，可重新命名精度。但初学者考试最好从一开始取 $\eps/2$，使结尾直接得到题目要求的 $<\eps$，不留下量词解释负担。

\Slide{58}{柯西推出收敛：先有界，再锁定子列极限}
\paragraph{先证明有界}在柯西条件中先取固定精度 $1$，得 $K_0$。固定 $m=K_0$，所有 $n\ge K_0$ 满足 $|x_n|\le|x_n-x_{K_0}|+|x_{K_0}|<1+|x_{K_0}|$。再与前面有限项的绝对值取最大值，便得到全列的统一界。
\paragraph{找到候选极限}由 Bolzano 定理，有子列 $x_{n_j}\to a$。
\paragraph{证明整列趋于它}任给 $\eps>0$，柯西条件给 $K$，使 $m,n\ge K$ 时距离小于 $\eps/2$。选一个足够大的 $j$，同时保证 $n_j\ge K$ 及 $|x_{n_j}-a|<\eps/2$。以后每个 $n\ge K$ 都满足
\[|x_n-a|\le|x_n-x_{n_j}|+|x_{n_j}-a|<\eps.\]
这里先选一个共同的子列项作为桥梁，而不是对每个 $n$ 随意更换未经控制的指标。

\Slide{59}{柯西充分性的第二证法：尾部下确界}
柯西列有界的证明同上一页。定义 $y_n=\inf\{x_k:k\ge n\}$。删掉第一项后尾集缩小，下确界只能上升：$y_{n+1}\ge y_n$。原数列的统一上下界也界住 $y_n$，所以 $y_n$ 单调有界，存在极限 $a$。
\paragraph{几何解释}$y_n$ 是从第 $n$ 项起所有数的“下边界”。它不一定是实际某一项，但总有 $y_n\le x_n$。柯西条件将迫使实际项与这条下边界越来越接近。
\paragraph{原页勘误}$|nn|$ 应为 $|x_n|$。不用把下确界视作最小值；即使尾集没有最小项，定义与证明仍然有效。

\Slide{60}{尾部下确界证法：修正指标后完成}
原页混用了取下确界的哑指标与自由指标。以下使用 $j$ 表示被取遍的指标，避免混淆。
给 $\eps>0$，取 $K_1$ 使 $n\ge K_1$ 时 $|y_n-a|<\eps/2$；再取 $K_2$ 使 $n,j\ge K_2$ 时 $|x_n-x_j|<\eps/2$。
固定 $n\ge\max(K_1,K_2)$。对\emph{所有} $j\ge n$，有 $x_j>x_n-\eps/2$，所以取下确界后
\[y_n\ge x_n-\eps/2,\qquad 0\le x_n-y_n\le\eps/2.\]
从而 $|x_n-a|\le|x_n-y_n|+|y_n-a|<\eps$，故 $x_n\to a$。
\paragraph{细节}取确界后严格不等号可能变成非严格号，不能保留原页所有 $<$。由于另一误差严格小于 $\eps/2$，最终仍得到严格 $<\eps$。

\Slide{61}{第三证法：尾部上确界}
定义 $z_n=\sup\{x_k:k\ge n\}$。尾集合变小使上确界不增，故 $z_n$ 递减；柯西列有界又保证 $z_n$ 有下界。因此单调收敛定理给 $z_n\to a$。
\paragraph{对照下确界方法}下包络 $y_n$ 从下方上升，上包络 $z_n$ 从上方下降；柯西条件意味着两者之间的尾部振幅趋零。两种方法都是把不单调的数列转为单调边界。
\paragraph{原页勘误}同样应把 $|nn|$ 改为 $|x_n|$。因为 $z_n$ 是上确界，所以 $x_n\le z_n$，不能反过来。下一页用这一方向完成估计。

\Slide{62}{尾部上确界证法：正确的误差估计}
任给 $\eps>0$，取 $K$ 同时保证 $n\ge K$ 时 $|z_n-a|<\eps/2$，且 $n,j\ge K$ 时 $|x_n-x_j|<\eps/2$。
固定 $n\ge K$，对所有 $j\ge n$ 有 $x_j<x_n+\eps/2$，所以
\[x_n\le z_n=\sup_{j\ge n}x_j\le x_n+\eps/2.\]
于是 $|x_n-a|\le|x_n-z_n|+|z_n-a|<\eps$，整列收敛。
\paragraph{课件指标修正}不能把 $\sup_{j\ge K}x_j$ 写成任意的 $z_{n_1}$，两者的尾集起点不同。本页证明固定 $n$ 后统一对 $j\ge n$ 取上确界，便消除了这个问题。上下包络证法任选一种熟练掌握即可，但应认识两种思路。

\Slide{63}{开覆盖与有限子覆盖}
开集在实数轴上的意思是每个点周围都有一小段区间仍留在集合内。开覆盖 $\mathcal H$ 是一族开集，使
\[S\subseteq\bigcup_{H\in\mathcal H}H.\]
\paragraph{原页勘误}覆盖公式必须有并集符号 $\bigcup$；不能把点的集合 $S$ 直接当成开集族的子集。
有限子覆盖要求从\emph{原来给定}的覆盖中选出 $H_1,\ldots,H_m$，仍有 $S\subseteq\bigcup_{i=1}^mH_i$，不能另造一个大开集来代替。
\paragraph{Heine--Borel}实数轴中每个闭有界集合都有该性质，称为紧。$(0,1)$ 被 $(1/n,1)$（$n\ge2$）覆盖，却没有有限子覆盖，因为有限多个的并仍漏掉靠近 $0$ 的点。

\Slide{64}{有限覆盖定理：把坏区间不断二分}
先证明闭区间 $[D,U]$。反设存在一个开覆盖，无法抽出有限子覆盖。二分成两个闭半区间，至少一个仍无法有限覆盖；否则两边各取有限个，合起来仍有限，便覆盖了整个区间，矛盾。
如此递归，得到每个都“无法有限覆盖”的闭区间 $I_n$，且嵌套、长度 $(U-D)/2^n\to0$。闭区间套定理给交点 $a$。
\paragraph{思路}我们暂时不直接挑覆盖集，而是把“失败的位置”压缩到一个点。开集在该点周围会留出余地，下一页正用这个余地制造矛盾。$D=U$ 的单点区间只需选一个覆盖它的开集，已经成立。

\Slide{65}{有限覆盖定理：矛盾与一般闭集情形}
既然原族覆盖 $[D,U]$，有 $H_0\in\mathcal H$ 含 $a$。$H_0$ 开，故有 $\delta>0$ 使 $(a-\delta,a+\delta)\subseteq H_0$。取 $n$ 使 $|I_n|<\delta$。由于 $a\in I_n$，整个 $I_n$ 都在该邻域内，于是仅一个 $H_0$ 就覆盖 $I_n$，与构造矛盾。
\paragraph{补足课件未写的一步}定理针对任意闭有界 $S$，而非只有区间。取 $S\subseteq[D,U]$。因为 $S$ 闭，$\R\setminus S$ 开，把它加入原开覆盖，便覆盖整个 $[D,U]$。用已证结论取有限子覆盖，再删除其中的 $\R\setminus S$（若有），剩下的原覆盖集仍覆盖 $S$。空集情形取空覆盖即可。
\paragraph{原页记号}邻域中心须统一写交点 $a$，不能在证明中突然换成未定义的 $x$。

\Slide{66}{完备性定理的第一张复习地图}
本页复习分割、确界、单调收敛和区间套。它们的联系比孤立背诵更重要：确界帮助确定边界，单调性让后续项留在边界附近，区间套把两个单调端点结合起来。
\paragraph{逐条检查条件}分割两侧必须非空；确界存在要非空并朝相应方向有界；单调收敛要单调与有界同时满足；区间套要非空、闭、有界、嵌套，长度趋零才有唯一性。
\paragraph{原页更精确的结论}区间套交集不仅包含 $[D,U]$，而是\emph{等于} $[D,U]$。闭卷训练：不看前文写出 $d_n\le u_m$ 对任意 $m,n$ 成立的理由，这一句串起整个区间套证明。

\Slide{67}{完备性定理的第二张复习地图}
Bolzano 定理从有界数列抽出收敛子列；Weierstrass 定理为有界无限集合找聚点；柯西准则把“尾项互近”转成收敛；Heine--Borel 把任意开覆盖压缩成有限个。
\paragraph{容易混淆的结论}有界数列不一定整列收敛，如 $(-1)^n$；它只是有收敛子列。有聚点不保证集合包含该点，如 $\{1/n:n\in\N\}$。闭性才保证包含所有聚点。
\paragraph{证明工具选择}看见“有界数列”，考虑抽子列；看见“无穷多个不同点”，考虑聚点；看见“任意两项”，考虑柯西；看见一族局部邻域，考虑有限覆盖。复习时各写一个反例，检查自己是否记住假设。